\chapter{Thigh Pain}

\section*{Definition}
Thigh pain refers to discomfort in the region between the hip and knee, arising from the muscles (quadriceps, hamstrings, adductors), bones (femur), nerves, or blood vessels. The etiology ranges from benign muscle strain to life-threatening conditions such as deep vein thrombosis (DVT) or necrotizing fasciitis.

\section*{Pathophysiology}
Muscle strains involve tearing of myofibrils at the myotendinous junction due to eccentric overload, with Grade I (microtears), Grade II (partial tear), or Grade III (complete rupture). Stress fractures of the femoral neck or shaft result from repetitive submaximal loading exceeding bone remodeling capacity. Radicular pain arises from compression of lumbar nerve roots (L1--L4) innervating the thigh, while vascular pain from DVT is caused by venous distension and perivascular inflammation.

\section*{Causes}
\begin{itemize}
  \item \textbf{Muscular:} Quadriceps or hamstring strain, adductor strain, contusion, myositis, tendinopathy
  \item \textbf{Skeletal:} Femoral stress fracture, pathologic fracture, osteoid osteoma, bone metastasis
  \item \textbf{Neurologic:} Lumbar radiculopathy (L2--L4), femoral neuropathy, meralgia paresthetica (lateral femoral cutaneous nerve), diabetic amyotrophy
  \item \textbf{Vascular:} Deep vein thrombosis, chronic venous insufficiency, peripheral artery disease
  \item \textbf{Joint-referred:} Hip osteoarthritis, hip labral tear, femoroacetabular impingement, sacroiliac joint pain
  \item \textbf{Infectious:} Cellulitis, pyomyositis, osteomyelitis, septic hip (referred)
  \item \textbf{Hematologic:} Hemarthrosis, sickle cell vaso-occlusive crisis
\end{itemize}

\section*{When to See a Doctor}
See your doctor if:
\begin{itemize}
  \item Acute onset of severe thigh pain with swelling, warmth, and redness (rule out DVT or infection)
  \item Pain with weight-bearing after trauma (rule out fracture)
  \item Thigh pain with neurological deficits (weakness, numbness, foot drop)
  \item Systemic symptoms: fever, night pain, unexplained weight loss
\end{itemize}

\section*{Related Symptoms}
\begin{itemize}
  \item Pain location: anterior (quadriceps), posterior (hamstrings), medial (adductors), lateral (trochanteric)
  \item Swelling, bruising, or palpable mass
  \item Weakness of knee extension (femoral nerve) or knee flexion (sciatic nerve)
  \item Antalgic gait or inability to bear weight
  \item Homan sign (calf pain with dorsiflexion) in DVT
\end{itemize}
\textbf{Important:} Acute anterior thigh pain with inability to bear weight after a sprinting injury suggests a quadriceps or rectus femoris rupture, often requiring surgical repair.

\section*{Self-Care / Home Management}
\begin{itemize}
  \item RICE (rest, ice, compression, elevation) for acute muscular injuries
  \item Over-the-counter NSAIDs or acetaminophen for mild to moderate pain
  \item Gradual return to activity after muscle strain (avoid early heavy loading)
  \item Stretching of the quadriceps, hamstrings, and hip flexors
  \item Use of crutches or cane if weight-bearing is painful
\end{itemize}

\section*{How Your Doctor Will Evaluate You}
\begin{itemize}
  \item Plain radiography of the femur and hip for fracture or bony lesion
  \item MRI for suspected stress fracture, occult fracture, muscle injury, or soft tissue mass
  \item Ultrasound for DVT (Doppler), collection of pus, or muscle tear evaluation
  \item Laboratory studies: CBC, ESR, CRP, CK, D-dimer if DVT suspected
  \item Venography or CT pulmonary angiography for confirmed thrombosis
  \item Nerve conduction studies for suspected femoral or lumbar radiculopathy
\end{itemize}

\section*{Treatment Options}
\begin{itemize}
  \item Immobilization and protected weight-bearing for fractures
  \item NSAIDs and physical therapy for muscle strains and tendinopathy
  \item Anticoagulation for DVT (direct oral anticoagulants or low-molecular-weight heparin)
  \item Antibiotics and surgical removal of dead tissue for pyomyositis or necrotizing fasciitis
  \item Image-guided corticosteroid injection for trochanteric bursitis
  \item Surgical fixation for femoral neck or shaft fractures
  \item Pain management and surgical surgical removal for bone tumors
\end{itemize}

\clearpage
